\documentclass[12pt]{article}
\usepackage[utf8]{inputenc}
\usepackage[T1]{fontenc}
\usepackage{hyperref}
\title{Silimo Language Resources for K-12 Curriculum Development}
\author{Compiled by Claude Sonnet 5.5 (claude-sonnet-5-5) for Andr\'as Kornai}
\date{}
\begin{document}
\maketitle

\section{Introduction}
Silimo (ISO 639-3 wul; also called Wulik or South Ngalik) is a Trans-New Guinea language of the highlands of Indonesian Papua, with about 5,000 speakers (a 1987 figure). It is distinct from Yali, whose varieties include Angguruk and North Ngalik. Resources for Yali do not apply to Silimo. Few resources turned up in the searches for this report.

\section{Primary Educational Resources}
No primer, textbook or school curriculum for Silimo turned up in the searches.

\section{Linguistic and Reference Materials}
A short encyclopedic overview is at \url{https://en.wikipedia.org/wiki/Silimo_language}. The searches turned up no Silimo dictionary, grammar or orthography description.

\section{Religious and Community Resources}
Scripture Earth lists Silimo scripture materials, including audio and text versions: \url{https://www.scriptureearth.org/00eng.php?iso=wul}. Joshua Project's page for South Ngalik in Indonesia is at \url{https://joshuaproject.net/people_groups/13332}. It reports that a full Bible translation was completed in 2023 after nearly five decades of work. I did not confirm that this is the same translation as the Scripture Earth listing.

\section{Digital and Audio Resources}
The Scripture Earth page mentions Bible.is audio and Global Recordings Network recordings. I did not see a direct URL for either, so they are not listed in the JSON.

\section{Curriculum Development Context}
The language is not known to be taught in schools, and little published material exists. A curriculum would need new work: orthography and literacy materials built from the Bible translation text, and community-led documentation.
\end{document}